% !PRINTED 44
\weberpage{44}
%<<CONT
\gesperrt{Abel's theorem}, and it is to be conceived and derived here in this way as well\footnote{This far-reaching theorem is found derived with magnificent simplicity in an article of barely two pages in the fourth volume of \gesperrt{Crelles} Journal ``D\'emonstration d'une propri\'et\'e g\'en\'erale d'une certaine classe de fonctions transcendentes''; Oeuvres compl\`etes de N. H. \gesperrt{Abel}. Nouvelle \'edition, T.\,I, p.\,515. A detailed presentation and applications are given in the great posthumous memoir: ``M\'emoire sur une propri\'et\'e g\'en\'erale d'une classe tr\`es \'etendue de fonctions transcendentes''.}. We preface the proof with the following elementary algebraic theorem:

\gesperrt{If $F(x)$ is an integral rational function of degree $n$ without multiple factors, $F'(x)$ its derivative, if further $x_1, x_2, \ldots, x_n$ are the roots of the equation $F(x) = 0$ and $\varphi(x)$ an integral function of $x$ whose degree is at most $= n-2$, then}
\weq{2}{\frac{\varphi(x_1)}{F'(x_1)} + \frac{\varphi(x_2)}{F'(x_2)} + \cdots \frac{\varphi(x_n)}{F'(x_n)} = 0.} % [sic] no plus sign is printed after the ellipsis

The proof of this theorem follows immediately from the decomposition of the rational fraction
\[
\frac{x\,\varphi(x)}{F(x)}
\]
into partial fractions
\[
\frac{x_1\varphi(x_1)}{F'(x_1)(x-x_1)} + \frac{x_2\varphi(x_2)}{F'(x_2)(x-x_2)} + \cdots + \frac{x_n\varphi(x_n)}{F'(x_n)(x-x_n)},
\]
if we there put $x = 0$ [Bd.\,I, \S\ 15 (9)].

Now let $P$, $Q$ denote integral functions of $x$ of the degrees $m$ and $m-2$, and we ask for the values of $x$, $\sqrt{f(x)}$, for which the function
\weq{3}{P + Q\sqrt{f(x)}}
vanishes. There are $2m$ such pairs of values (points), and one finds the values of $x$ as the roots of the equation of degree $2m$:
\weq{4}{F(x) = P^2 - Q^2 f(x) = 0,}
and the corresponding values of $\sqrt{f(x)}$ from
\weq{5}{P + Q\sqrt{f(x)} = 0.}

We now assume that the coefficients in the functions $P$, $Q$ are variable, either independent variables or dependent in some way on other variables.
